\documentclass[letter, 10pt]{article}
\usepackage[utf8]{inputenc}
\usepackage[T1]{fontenc}
\usepackage[spanish, es-nodecimaldot, es-tabla, es-noquoting]{babel}
\usepackage{amsfonts}
\usepackage{amsmath}
\usepackage{amssymb}
\usepackage[dvips]{graphicx}
\usepackage{url}
\usepackage{booktabs}
\usepackage{tikz}
\usetikzlibrary{babel}
\usepackage[top=3cm,bottom=3cm,left=3.5cm,right=3.5cm,footskip=1.5cm,headheight=1.5cm,headsep=.5cm,textheight=3cm]{geometry}

\newcommand{\R}{\mathbb{R}}

\begin{document}
\title{Inteligencia Artificial \\ \begin{Large}Estado del Arte: Problema de Programación de Conductores de Buses (\emph{Bus Driver Scheduling Problem}, BDSP)\end{Large}}
\author{Vicente Luongo \\ ROL: 202073637-5}
\date{\today}
\maketitle


%--------------------No borrar esta secci\'on--------------------------------%
\section*{Evaluaci\'on}

\begin{tabular}{ll}
Resumen (5\%): & \underline{\hspace{2cm}} \\
Introducci\'on (5\%):  & \underline{\hspace{2cm}} \\
Definici\'on del Problema (10\%):  & \underline{\hspace{2cm}} \\
Estado del Arte (40\%):  & \underline{\hspace{2cm}} \\
Modelo Matem\'atico (20\%): &  \underline{\hspace{2cm}}\\
Conclusiones (15\%): &  \underline{\hspace{2cm}}\\
Bibliograf\'ia (5\%): & \underline{\hspace{2cm}}\\
 &  \\
\textbf{Nota Final (100\%)}:   & \underline{\hspace{2cm}}
\end{tabular}
%---------------------------------------------------------------------------%
\vspace{2cm}


\begin{abstract}
El \emph{Bus Driver Scheduling Problem} (BDSP) consiste en construir los turnos diarios de trabajo de los conductores de una empresa de transporte público, de modo que cada tramo de recorrido de los buses sea cubierto exactamente por un conductor, respetando la legislación laboral y los convenios colectivos. Se trata de un problema de optimización combinatoria NP-duro, cuyo costo impacta directamente en la operación de las empresas, dado que la mano de obra constituye una de sus principales partidas de gasto. En este documento se define el problema y la variante a estudiar, cuyo objetivo es minimizar el tiempo pagado total bajo reglas de descanso, sobretiempo, trabajo nocturno, tiempo mínimo pagado y duración máxima del turno. Se revisa la literatura desde los modelos clásicos de partición y cobertura de conjuntos hasta las metaheurísticas y los enfoques híbridos recientes, y se propone un modelo de programación lineal entera mixta que captura todas las reglas de la variante estudiada.
\end{abstract}

\section{Introducci\'on}
La planificación de un sistema de transporte público por buses se descompone habitualmente en una secuencia de subproblemas: diseño de la red, construcción de itinerarios (\emph{timetabling}), programación de vehículos (\emph{vehicle scheduling}), programación de tripulaciones (\emph{crew scheduling}) y asignación de turnos a personas concretas en un horizonte de semanas (\emph{crew rostering})~\cite{IbarraRojas2015,DeLeone2011a}. El BDSP corresponde a la etapa de programación de tripulaciones: dado el plan de vehículos ya fijado, se deben formar los turnos (\emph{duties}) que cubren todo el trabajo de conducción de un día.

La motivación del problema es principalmente económica y legal. El costo de los conductores es uno de los componentes más importantes del costo operacional de una empresa de transporte~\cite{WrenRousseau1995}, y una asignación deficiente genera horas extraordinarias, tiempos muertos pagados y un número innecesariamente alto de conductores. Al mismo tiempo, los turnos deben cumplir reglas laborales estrictas (descansos mínimos, jornada máxima, recargos por trabajo nocturno), lo que hace que la construcción manual de turnos sea lenta y difícilmente óptima. Adicionalmente, el problema es NP-duro incluso en versiones muy simplificadas~\cite{Fischetti1987,Fischetti1989}, lo que justifica el estudio de técnicas de inteligencia artificial y de investigación de operaciones para resolverlo.

El propósito de este informe es presentar el estado del arte del BDSP y formular un modelo matemático para la variante asignada en el curso, basada en las reglas de una empresa brasileña descritas por Constantino et al.~\cite{Constantino2017}. El documento se organiza como sigue: la Sección~\ref{sec:def} define el problema, su terminología, la variante estudiada y los problemas relacionados; la Sección~\ref{sec:eda} revisa los principales métodos de resolución propuestos en la literatura; la Sección~\ref{sec:modelo} presenta dos modelos matemáticos, uno general y otro específico para la variante estudiada; y la Sección~\ref{sec:conc} expone las conclusiones del estudio.

\section{Definici\'on del Problema}\label{sec:def}

\subsection{Terminología}
El BDSP recibe como entrada la programación de vehículos, es decir, la secuencia de viajes que realiza cada bus durante el día (su \emph{bloque}). A lo largo de cada bloque existen \emph{puntos de relevo} (\emph{relief opportunities}): pares lugar--instante en que un conductor puede dejar el vehículo y otro tomarlo. El trabajo comprendido entre dos puntos de relevo consecutivos se denomina \emph{tarea} o \emph{tramo}, y es la unidad mínima e indivisible de trabajo que se asigna a un conductor~\cite{Constantino2017}. Las paradas intermedias de un tramo no son relevantes, pues solo afectan su duración.

Un \emph{turno} (\emph{duty}) es la secuencia ordenada de tramos que realiza un conductor en un día. Se denomina \emph{spreadover} al tiempo transcurrido entre el inicio del primer tramo y el término del último. Entre dos tramos consecutivos de un turno puede existir una pausa, que puede ser un \emph{descanso} (no pagado, típicamente destinado a alimentación) o un \emph{tiempo ocioso} (pagado). Una secuencia de tramos trabajados sin un descanso intermedio se denomina \emph{stretch}~\cite{Constantino2017,DeLeone2011a}.

\subsection{Descripción general}
En términos generales, el BDSP busca un conjunto de turnos tal que:
\begin{itemize}
\item \textbf{Cobertura:} cada tramo es realizado por exactamente un conductor; no se permite que dos conductores compartan un tramo ni que un tramo quede sin cubrir.
\item \textbf{Continuidad temporal y espacial:} un conductor solo puede comenzar un tramo si ya terminó el anterior y si se encuentra en el lugar donde este comienza. Sí puede cambiar de bus o de línea en cualquier punto de relevo.
\item \textbf{Factibilidad laboral:} cada turno cumple las reglas de la legislación y de los convenios (descansos, jornada, sobretiempo, etc.).
\item \textbf{Optimalidad:} se minimiza el costo total, usualmente medido como el tiempo pagado a los conductores, el número de conductores o una combinación de ambos.
\end{itemize}
Las \emph{variables} del problema corresponden, por lo tanto, a qué conductor realiza cada tramo (lo que determina implícitamente la secuencia de cada turno), y las decisiones derivadas de ello: cuál pausa se considera descanso, cuánto sobretiempo se genera, etc.

\subsection{Variante estudiada}\label{sec:variante}
La variante asignada corresponde a la estudiada por Constantino et al.~\cite{Constantino2017}, con reglas provenientes de la legislación laboral brasileña y los convenios sindicales de una empresa de transporte urbano. Cada tramo queda descrito por su línea, punto de partida, punto de término, hora de partida, duración y bus. Las reglas, parametrizadas en cada instancia, son las siguientes (entre paréntesis, los valores del caso de referencia):
\begin{enumerate}
\item \textbf{Stretch máximo y descanso mínimo.} Un conductor no puede trabajar más de $T^{\max}$ (360 min) sin un descanso; si el turno excede ese tiempo, debe incluir un descanso de al menos $R^{\min}$ (90 min).
\item \textbf{Umbral de descanso pagado.} El descanso no se paga hasta un umbral $R^{\text{pag}}$ (300 min); el exceso por sobre él se considera tiempo ocioso y se paga.
\item \textbf{Tiempo pagado mínimo.} Todo turno se paga por al menos $P^{\min}$ (432 min), aunque el trabajo efectivo sea menor.
\item \textbf{Sobretiempo.} El tiempo trabajado por sobre la jornada normal es sobretiempo, no puede superar $O^{\max}$ (120 min) y recibe un recargo de $\beta_O$ (50\%).
\item \textbf{Trabajo nocturno.} La conducción entre las 22:00 y las 05:00 se corrige por un factor $f_N$ (60/52, es decir, 52 minutos trabajados se cuentan como 60) y recibe un recargo de $\beta_N$ (20\%).
\item \textbf{Duración máxima del turno.} El spreadover, incluyendo descansos no pagados, no puede superar $S^{\max}$ (780 min).
\item \textbf{Cambio de línea.} Los cambios de línea dentro de un turno están permitidos, pero se penalizan con $p_L$ (1 min) cada uno.
\end{enumerate}
El objetivo es minimizar el tiempo pagado total de todos los conductores más las penalizaciones por cambio de línea. Nótese que, a diferencia de otras variantes, el número de conductores no se minimiza explícitamente; sin embargo, la regla del tiempo pagado mínimo induce a usar pocos turnos bien aprovechados.

La Figura~\ref{fig:ejemplo} muestra la instancia de ejemplo entregada en el curso, con seis tramos de dos líneas. Un único conductor no podría cubrirlos todos, pues el spreadover resultante (05:50--20:35, 885 min) excede $S^{\max}$. La solución natural usa dos conductores, uno por bus, cada uno con menos de 200 minutos de trabajo; por la regla de tiempo mínimo, cada uno se paga 432 minutos, lo que da un costo de 864. El ejemplo ilustra cómo el tiempo mínimo pagado y el tiempo ocioso dominan el costo cuando los turnos son cortos.

\begin{figure}[ht]
\centering
\begin{tikzpicture}[x=0.58cm,y=1cm,font=\scriptsize]
  % franjas nocturnas
  \fill[gray!20] (0,-0.4) rectangle (1,2.0);
  \fill[gray!20] (18,-0.4) rectangle (19,2.0);
  \node[gray] at (0.5,2.2) {noche};
  \node[gray] at (18.5,2.2) {noche};
  % eje de tiempo (horas 4 a 23)
  \draw[->] (0,-0.4) -- (19.3,-0.4);
  \foreach \h in {4,6,...,22} {
    \draw (\h-4,-0.45) -- (\h-4,-0.35);
    \node[below] at (\h-4,-0.45) {\h:00};
  }
  % filas de buses
  \node[left] at (0,1.3) {Bus 242 (línea 202)};
  \node[left] at (0,0.4) {Bus 98 (línea 235)};
  % tramos bus 242: 5 (05:50-06:00), 6 (06:00-07:20)
  \draw[fill=teal!40] (1.833,1.1) rectangle (2.0,1.5);
  \draw[fill=teal!40] (2.0,1.1) rectangle (3.333,1.5);
  \node[above] at (1.8,1.5) {5};
  \node[above] at (2.67,1.5) {6};
  \node[below] at (1.833,1.1) {G};
  \node[below] at (2.35,1.1) {T};
  \node[below] at (3.333,1.1) {T};
  % tramos bus 98: 1 (17:20-17:35), 2 (17:35-18:40), 3 (18:40-19:50), 4 (19:50-20:35)
  \draw[fill=orange!40] (13.333,0.2) rectangle (13.583,0.6);
  \draw[fill=orange!40] (13.583,0.2) rectangle (14.667,0.6);
  \draw[fill=orange!40] (14.667,0.2) rectangle (15.833,0.6);
  \draw[fill=orange!40] (15.833,0.2) rectangle (16.583,0.6);
  \node[above] at (13.3,0.6) {1};
  \node[above] at (14.12,0.6) {2};
  \node[above] at (15.25,0.6) {3};
  \node[above] at (16.2,0.6) {4};
  \node[below] at (13.2,0.2) {G};
  \node[below] at (13.75,0.2) {T};
  \node[below] at (14.667,0.2) {T};
  \node[below] at (15.833,0.2) {T};
  \node[below] at (16.583,0.2) {G};
\end{tikzpicture}
\caption{Instancia de ejemplo del curso: seis tramos sobre dos buses. Las letras indican los puntos de relevo (G: garaje, T: terminal) y las zonas grises, el horario nocturno.}
\label{fig:ejemplo}
\end{figure}

\subsection{Problemas relacionados y variantes}
El BDSP pertenece a la familia de problemas de \emph{crew scheduling}, que también incluye la programación de tripulaciones aéreas y ferroviarias; en todos ellos se deben cubrir tareas con secuencias factibles de trabajo, aunque las reglas y la escala difieren. Está estrechamente ligado al \emph{vehicle scheduling}, que determina los bloques, y al \emph{crew rostering}, que encadena los turnos diarios en calendarios semanales o mensuales considerando aspectos de calidad de vida~\cite{IbarraRojas2015}. Desde el punto de vista teórico, se relaciona con el \emph{fixed job scheduling} con restricciones de spreadover y de tiempo de trabajo, cuya NP-dureza fue demostrada por Fischetti et al.~\cite{Fischetti1987,Fischetti1989}.

Las variantes más conocidas del BDSP se diferencian en: (i) el \emph{objetivo}, que puede ser minimizar el número de conductores (versión \emph{unicost})~\cite{Yunes2005}, el costo total~\cite{Constantino2017}, o ambos, incluso de forma multiobjetivo~\cite{Lourenco2001}; (ii) el \emph{alcance}, pues algunas versiones resuelven la programación de vehículos y conductores de forma integrada~\cite{Huisman2005}; (iii) la existencia de uno o varios \emph{depósitos}, con la exigencia de terminar el turno en el depósito de partida~\cite{DeLeone2011a}; (iv) la separación por línea frente a un enfoque de red, en que un conductor puede cambiar de línea~\cite{Portugal2009,Constantino2017}; y (v) la complejidad de las reglas de pausas, como en el caso austriaco estudiado por Kletzander y Musliu~\cite{Kletzander2020}.

\section{Estado del Arte}\label{sec:eda}

\subsection{Origen y primeros enfoques}
El BDSP se estudia desde la década de 1960, inicialmente mediante heurísticas que imitaban el proceso manual de los planificadores~\cite{WrenRousseau1995}. La revisión de Bodin et al.~\cite{Bodin1983} ya presentaba en 1983 un panorama amplio de la programación de vehículos y tripulaciones, y los avances posteriores se han difundido principalmente en la serie de conferencias \emph{Computer-Aided Scheduling of Public Transport} (hoy CASPT)~\cite{Constantino2017}. Wren y Rousseau~\cite{WrenRousseau1995} resumen la evolución del área y consolidan el enfoque que dominaría las décadas siguientes: generar un gran conjunto de turnos factibles y seleccionar entre ellos mediante programación entera.

\subsection{Métodos exactos basados en partición y cobertura de conjuntos}
La formulación más utilizada es la de \emph{set partitioning} (SPP) o \emph{set covering} (SCP): cada columna de la matriz representa un turno factible y cada fila un tramo; se busca el subconjunto de columnas de costo mínimo que cubre cada fila exactamente una vez (SPP) o al menos una vez (SCP)~\cite{Portugal2009,DeLeone2011a}. Su principal ventaja es que toda la complejidad de las reglas laborales queda encapsulada en la generación de columnas, por lo que el modelo es independiente de la legislación. Su principal desventaja es que el número de turnos posibles crece de manera exponencial con el número de tramos.

Para enfrentar esta dificultad se han seguido dos caminos. El primero es la enumeración heurística de un subconjunto de turnos ``buenos'', como en el sistema TRACS~II, que combina programación entera con heurísticas de reducción y generación de columnas en el nodo raíz~\cite{Fores2002}. El segundo es la \emph{generación de columnas} (CG), introducida para este problema por Desrochers y Soumis~\cite{DesrochersSoumis1989}, donde un problema maestro (SPP/SCP relajado) se resuelve junto con un subproblema de camino más corto con restricciones de recursos que genera nuevos turnos de costo reducido negativo. Yunes et al.~\cite{Yunes2005} combinaron CG con programación por restricciones para resolver instancias reales de hasta 246 tareas minimizando el número de conductores, y Chen y Shen~\cite{ChenShen2013} propusieron mejoras de CG para instancias de hasta 701 viajes. Curtis et al.~\cite{Curtis1999} emplearon programación por restricciones guiada por la relajación lineal, lo que permite incorporar reglas no lineales. En el ámbito integrado, Huisman et al.~\cite{Huisman2005} resolvieron simultáneamente vehículos y conductores con múltiples depósitos mediante CG y relajación lagrangiana.

Como alternativa a SPP/SCP, De Leone et al.~\cite{DeLeone2011a} propusieron un modelo compacto de programación entera en que las variables binarias indican qué tramo ocupa cada posición de cada turno, con restricciones explícitas de spreadover, tiempo de trabajo, cantidad de tiempos ociosos y retorno al depósito. Resuelto con CPLEX, el modelo encontró óptimos en instancias de hasta 25 tramos, pero en 70 tramos solo consiguió una primera solución factible para uno de los tres objetivos, lo que evidencia la limitación de los métodos exactos directos.

La línea exacta más reciente es la de Kletzander et al., que aplican \emph{branch and price} a un BDSP con reglas de pausa austriacas muy complejas~\cite{Kletzander2021}, y luego lo integran con búsqueda en vecindades grandes (LNS), usando CG o \emph{branch and price} para reparar subproblemas y reutilizando las columnas generadas~\cite{Kletzander2026}.

\subsection{Metaheurísticas y heurísticas}
Cuando las reglas son difíciles de expresar linealmente o las instancias son muy grandes, las metaheurísticas han sido la opción preferida. Las principales familias son:

\textbf{Algoritmos genéticos.} Wren y Wren~\cite{WrenWren1995} propusieron uno de los primeros algoritmos genéticos para el problema, aplicado sobre la formulación de cobertura de conjuntos. Dias et al.~\cite{Dias2002} extendieron la formulación SPP/SCP con un algoritmo genético capaz de considerar varios criterios de calidad simultáneamente, y Li y Kwan~\cite{LiKwan2003} incorporaron lógica difusa para evaluar la calidad de los turnos dentro de un algoritmo genético.

\textbf{Búsqueda tabú.} Shen y Kwan~\cite{ShenKwan2001} desarrollaron HACS, una búsqueda tabú que opera directamente sobre los turnos, sin recurrir a la enumeración de columnas propia de SPP/SCP.

\textbf{GRASP.} Lourenço et al.~\cite{Lourenco2001} aplicaron GRASP~\cite{FeoResende1995}, junto con búsqueda tabú y algoritmos genéticos multiobjetivo, con una construcción voraz basada en el costo de los turnos y una búsqueda local de intercambio simple; estas técnicas se incorporaron al sistema comercial GIST. De Leone et al.~\cite{DeLeone2011a} construyen cada turno agregando tramos compatibles desde una lista restringida de candidatos (RCL) ordenada por el tiempo de espera, y mejoran la solución con tres vecindarios: \emph{1-swap} (intercambio de dos tramos entre turnos), \emph{variable neighborhood swap} (intercambio de sufijos de turnos, que permite fusionar turnos y reducir su número) y una generalización a $\rho$ puntos de corte. En instancias reales de hasta 161 viajes, GRASP obtuvo en segundos soluciones comparables e incluso mejores que un método exacto basado en SCP con columnas truncadas, que requería desde minutos hasta horas. Los mismos autores propusieron variantes multiarranque aleatorizadas~\cite{DeLeone2011b} y la hibridación de GRASP con \emph{rollout}~\cite{DAnnibale2007}.

\textbf{Búsqueda en vecindario variable.} Ma et al.~\cite{Ma2016} emplearon VNS en un caso real de Beijing con 501 tareas.

\textbf{Heurística de múltiples asignaciones.} Constantino et al.~\cite{Constantino2017} proponen GraphBDSP, un algoritmo determinista que representa los tramos como vértices de un grafo multipartito dirigido y acíclico, donde los arcos indican que un tramo puede seguir a otro. Los vértices se agrupan en capas de tramos mutuamente incompatibles, y cada turno es un camino en el grafo. En la fase constructiva, recorre las capas y en cada una resuelve un problema de asignación lineal entre los turnos parciales y los tramos de la capa. En la fase de mejora, corta todos los turnos en una misma capa y recombina las partes izquierdas con las derechas mediante un nuevo problema de asignación (lo que nunca empeora la solución), además de un procedimiento que intenta eliminar el sobretiempo de los turnos largos. La factibilidad se maneja asignando costo infinito a las combinaciones que violan alguna regla. Con un solo parámetro, resolvió una instancia real de 2.313 tareas en pocos minutos, obteniendo mejores costos que un SCP con columnas generadas heurísticamente. Además, los autores encontraron que guiar la búsqueda por tiempo ocioso y sobretiempo da mejores resultados que usar directamente el tiempo pagado.

\textbf{Enfoques recientes.} Kletzander y Musliu~\cite{Kletzander2020} formalizaron las reglas austriacas, publicaron un conjunto de instancias y propusieron una metaheurística de búsqueda local que mejoró significativamente las soluciones reales. También se han explorado hiperheurísticas basadas en CG~\cite{Li2015}.

\subsection{Síntesis comparativa}
La Tabla~\ref{tab:resumen} resume los trabajos más representativos. Se observa un crecimiento sostenido en el tamaño de las instancias y una convergencia hacia métodos híbridos.

\begin{table}[ht]
\centering
\small
\caption{Resumen de enfoques representativos para el BDSP (tamaños reportados por los autores o en~\cite{Constantino2017}).}
\label{tab:resumen}
\setlength{\tabcolsep}{4pt}
\begin{tabular}{@{}llll@{}}
\toprule
Referencia & Objetivo & Técnica & Instancia mayor \\
\midrule
Wren y Wren 1995~\cite{WrenWren1995} & Costo/conductores & Alg. genético & -- \\
Lourenço et al. 2001~\cite{Lourenco2001} & Multiobjetivo & GRASP, tabú, AG & Casos reales \\
Shen y Kwan 2001~\cite{ShenKwan2001} & Costo & Búsqueda tabú & -- \\
Fores et al. 2002~\cite{Fores2002} & Costo & SCP + CG (TRACS II) & Casos reales \\
Yunes et al. 2005~\cite{Yunes2005} & Conductores & CG + PR & 246 tareas \\
Huisman et al. 2005~\cite{Huisman2005} & Costo & CG + lagrangiano & 653 viajes \\
De Leone et al. 2011~\cite{DeLeone2011a} & Costo/conductores & MIP, GRASP & 161 viajes \\
Ma et al. 2016~\cite{Ma2016} & Costo & VNS & 501 tareas \\
Constantino et al. 2017~\cite{Constantino2017} & Tiempo pagado & Asignaciones sucesivas & 2.313 tareas \\
Kletzander et al. 2021--2026~\cite{Kletzander2021,Kletzander2026} & Costo & Branch and price + LNS & Casos reales \\
\bottomrule
\end{tabular}
\end{table}

En cuanto a la \textbf{representación}, conviven tres enfoques: (i) conjuntos de turnos preenumerados (columnas), propio de SPP/SCP; (ii) turnos como secuencias o caminos en un grafo de compatibilidad, usado por las metaheurísticas de búsqueda local y por GraphBDSP; y (iii) asignación directa tramo--conductor, donde el orden dentro del turno queda determinado por el horario. En cuanto a los \textbf{movimientos}, destacan el intercambio de tramos entre turnos, la reubicación de un tramo, el intercambio de sufijos (que permite fusionar o dividir turnos) y la recombinación masiva mediante problemas de asignación. Finalmente, el \textbf{manejo de restricciones} se realiza mayoritariamente descartando las soluciones infactibles (costo infinito o filtrado en la generación de columnas), y en menor medida mediante penalizaciones en la función objetivo.

\section{Modelo Matem\'atico}\label{sec:modelo}
Se presentan dos modelos. El primero es la formulación clásica de partición de conjuntos, útil como marco general. El segundo es un modelo compacto de programación lineal entera mixta (MILP) propuesto para la variante estudiada; toma como base las reglas y la función de costo de Constantino et al.~\cite{Constantino2017}, quienes indican que no conocen una formulación lineal para ellas, y adapta ideas del modelo de De Leone et al.~\cite{DeLeone2011a}, reemplazando las variables por posición por variables de arco.

\subsection{Modelo 1: partición de conjuntos}
Sea $J$ el conjunto de tramos y $\mathcal{P}$ el conjunto de todos los turnos factibles, es decir, aquellos que cumplen todas las reglas de la Sección~\ref{sec:variante}. Para cada $p\in\mathcal{P}$, $c_p$ es su costo (tiempo pagado más penalizaciones) y $a_{jp}=1$ si el turno $p$ contiene el tramo $j$. Con $z_p\in\{0,1\}$ igual a 1 si se usa el turno $p$:
\begin{align}
\min \; & \sum_{p\in\mathcal{P}} c_p\, z_p \\
\text{s.a.}\; & \sum_{p\in\mathcal{P}} a_{jp}\, z_p = 1 && \forall j\in J \label{eq:spp}\\
& z_p\in\{0,1\} && \forall p\in\mathcal{P}
\end{align}
La restricción~\eqref{eq:spp} garantiza que cada tramo sea cubierto exactamente una vez. Al reemplazar ``$=$'' por ``$\geq$'' se obtiene el SCP. El espacio de búsqueda es $2^{|\mathcal{P}|}$, y $|\mathcal{P}|$ crece exponencialmente con $|J|$, lo que motiva el uso de generación de columnas o de enumeración parcial.

\subsection{Modelo 2: MILP compacto para la variante estudiada}

\paragraph{Conjuntos y parámetros.}
\begin{itemize}
\item $J=\{1,\dots,n\}$: tramos. Para cada $j\in J$: instante de inicio $s_j$ y término $e_j=s_j+d_j$ (minutos desde las 00:00, con $d_j$ la duración), lugar de partida $a_j$ y de término $b_j$, línea $\ell_j$, y minutos de conducción nocturna $\nu_j=|[s_j,e_j]\cap\mathcal{N}|$, donde $\mathcal{N}$ es el horario entre las 22:00 y las 05:00.
\item $K=\{1,\dots,m\}$: conductores disponibles ($m\leq n$ es una cota superior, por ejemplo $m=n$).
\item $A=\{(i,j)\in J\times J : e_i\leq s_j,\ b_i=a_j\}$: pares compatibles, en que el conductor que termina $i$ puede iniciar $j$. Se agregan un nodo origen $0$ y un sumidero $n{+}1$, con arcos $(0,j)$ y $(j,n{+}1)$ para todo $j$. Como $e_i\leq s_j$ y $d_j>0$, el grafo es acíclico.
\item $g_{ij}=s_j-e_i$: pausa entre $i$ y $j$. $A^R=\{(i,j)\in A: g_{ij}\geq R^{\min}\}$: pausas que pueden ser descanso. $\bar g_{ij}=\min\{g_{ij},R^{\text{pag}}\}$: porción no pagada si la pausa es el descanso.
\item $\lambda_{ij}=1$ si $\ell_i\neq\ell_j$ (cambio de línea) y 0 en otro caso.
\item $T^{\max}, R^{\min}, R^{\text{pag}}, P^{\min}, O^{\max}, S^{\max}, f_N, \beta_N, \beta_O, p_L$: parámetros de la Sección~\ref{sec:variante}. $M$: constante suficientemente grande (por ejemplo, el largo del horizonte en minutos).
\end{itemize}

\paragraph{Variables.}
\begin{itemize}
\item $x_{jk}\in\{0,1\}$: 1 si el tramo $j$ es asignado al conductor $k$.
\item $y_{ijk}\in\{0,1\}$: 1 si el conductor $k$ realiza el tramo $j$ inmediatamente después de $i$ (incluye arcos desde $0$ y hacia $n{+}1$).
\item $r_{ijk}\in\{0,1\}$, $(i,j)\in A^R$: 1 si la pausa entre $i$ y $j$ es el descanso del conductor $k$.
\item $u_k\in\{0,1\}$: 1 si el conductor $k$ es utilizado.
\item $S_k, E_k\geq 0$: inicio y término del turno $k$.
\item $W_k\geq 0$: tiempo trabajado del turno $k$, corregido por el factor nocturno.
\item $O_k\geq 0$: sobretiempo del turno $k$.
\end{itemize}
Como supuestos de modelamiento se considera, en concordancia con~\cite{Constantino2017}, que cada turno tiene a lo sumo un descanso no pagado, que las demás pausas son tiempo ocioso pagado, y que no se exige volver al punto de partida al final del turno.

\paragraph{Función objetivo.}
\begin{equation}
\min\; \sum_{k\in K}\Big( P^{\min}u_k + (1+\beta_O)\,O_k + \beta_N f_N \sum_{j\in J}\nu_j x_{jk}\Big) + p_L\sum_{k\in K}\sum_{(i,j)\in A}\lambda_{ij}\,y_{ijk}
\label{eq:fo}
\end{equation}
El primer término es el tiempo pagado. Si el tiempo trabajado $W_k$ no supera $P^{\min}$, se paga $P^{\min}$ (la diferencia es tiempo ocioso); si lo supera, se paga $P^{\min}+O_k=W_k$ más el recargo $\beta_O O_k$ del sobretiempo. A ello se suma el recargo nocturno sobre los minutos nocturnos corregidos. El segundo término penaliza los cambios de línea.

\paragraph{Restricciones.}
\begin{align}
& \sum_{k\in K} x_{jk} = 1 && \forall j\in J \label{r:cob}\\
& \sum_{i:(i,j)\in A\cup\{(0,j)\}} y_{ijk} = x_{jk},\quad \sum_{l:(j,l)\in A\cup\{(j,n+1)\}} y_{jlk} = x_{jk} && \forall j\in J,\ k\in K \label{r:flujo}\\
& \sum_{j\in J} y_{0jk} = u_k,\quad \sum_{j\in J} y_{j,n+1,k} = u_k && \forall k\in K \label{r:uso}\\
& S_k = \sum_{j\in J} s_j\, y_{0jk},\quad E_k=\sum_{j\in J} e_j\, y_{j,n+1,k} && \forall k\in K \label{r:ini}\\
& E_k - S_k \leq S^{\max} && \forall k\in K \label{r:spread}\\
& r_{ijk}\leq y_{ijk} && \forall (i,j)\in A^R,\ k\in K \label{r:rest1}\\
& \sum_{(i,j)\in A^R} r_{ijk}\leq 1 && \forall k\in K \label{r:rest2}\\
& E_k - S_k \leq T^{\max} + M\sum_{(i,j)\in A^R} r_{ijk} && \forall k\in K \label{r:sinrest}\\
& e_i - S_k \leq T^{\max} + M(1-r_{ijk}),\quad E_k - s_j \leq T^{\max} + M(1-r_{ijk}) && \forall (i,j)\in A^R,\ k\in K \label{r:stretch}\\
& W_k = E_k - S_k - \sum_{(i,j)\in A^R}\bar g_{ij}\, r_{ijk} + (f_N-1)\sum_{j\in J}\nu_j x_{jk} && \forall k\in K \label{r:trab}\\
& O_k \geq W_k - P^{\min}u_k && \forall k\in K \label{r:ot}\\
& O_k \leq O^{\max} && \forall k\in K \label{r:otmax}\\
& u_k \geq u_{k+1} && k=1,\dots,m-1 \label{r:sim}
\end{align}

\paragraph{Explicación.}
\eqref{r:cob} asegura que cada tramo sea realizado por exactamente un conductor, de modo que ningún tramo quede sin cubrir ni sea compartido. \eqref{r:flujo} enlaza la asignación con la secuencia: cada tramo asignado a $k$ tiene exactamente un antecesor y un sucesor en el turno de $k$ (posiblemente el origen o el sumidero). Como los arcos solo existen entre tramos compatibles, esto impone que un conductor no inicie un tramo antes de terminar el anterior ni en un lugar distinto del que se encuentra. \eqref{r:uso} hace que cada conductor utilizado tenga un único camino desde el origen al sumidero; al ser el grafo acíclico, no se requieren restricciones de eliminación de subtours. \eqref{r:ini} define el inicio y el término del turno, que valen 0 para conductores no utilizados, y \eqref{r:spread} limita la duración total del turno (regla 6). \eqref{r:rest1} y \eqref{r:rest2} permiten designar como descanso a lo más una pausa del turno, de largo al menos $R^{\min}$. \eqref{r:sinrest} y \eqref{r:stretch} modelan la regla 1: si el turno no tiene descanso, su duración total no puede exceder $T^{\max}$; si lo tiene, tanto el tramo de trabajo anterior como el posterior al descanso deben durar a lo más $T^{\max}$. \eqref{r:trab} calcula el tiempo trabajado como el spreadover menos la parte no pagada del descanso (regla 2), considerando que los minutos nocturnos valen $f_N$ veces (regla 5). \eqref{r:ot} define el sobretiempo como el exceso sobre la jornada; al minimizar, $O_k=\max\{0,W_k-P^{\min}\}$ y así el primer término de~\eqref{eq:fo} implementa la regla 3. \eqref{r:otmax} limita el sobretiempo (regla 4). Por último, \eqref{r:sim} elimina simetrías entre conductores equivalentes, obligando a usar primero los de menor índice. Los dominios de las variables son los indicados anteriormente.

\paragraph{Espacio de búsqueda.}
Dado que dentro de un turno el orden de los tramos queda determinado por sus horarios, una solución queda completamente definida por las variables $x_{jk}$; las variables $y$, $r$ y las continuas se derivan de ellas. Por lo tanto, el espacio de búsqueda tiene tamaño $|\Omega|\leq m^{n}$ (o a lo más el número de Bell $B_n$ si se ignora la identidad de los conductores), del cual solo una fracción pequeña es factible. El modelo tiene $O(|A|\,m)$ variables binarias y restricciones, lo que, al igual que el modelo de De Leone et al.~\cite{DeLeone2011a}, lo hace práctico solo para instancias pequeñas. Sin embargo, sirve para definir de manera precisa la función de evaluación y las restricciones que deberá manejar un algoritmo de búsqueda.

\section{Conclusiones}\label{sec:conc}
El estudio muestra que el BDSP no es un problema único, sino una familia de problemas cuyas diferencias provienen principalmente de las reglas laborales de cada país y empresa y del objetivo considerado (costo, número de conductores o ambos). Por ello, las técnicas no resuelven exactamente el mismo problema, y sus resultados son difíciles de comparar: muchos trabajos usan instancias privadas, y solo recientemente se han publicado conjuntos de instancias de referencia~\cite{Constantino2017,Kletzander2020}.

Las técnicas se agrupan en dos grandes enfoques. Los métodos basados en SPP/SCP y generación de columnas separan limpiamente las reglas (en la generación de turnos) de la optimización y entregan cotas de calidad, pero su costo computacional crece rápidamente, y cuando las columnas se generan de forma heurística pierden la garantía de optimalidad: en~\cite{Constantino2017}, el SCP con columnas truncadas fue superado por una heurística. Las metaheurísticas (GRASP, tabú, genéticos, VNS) y las heurísticas especializadas escalan a instancias de miles de tramos, pero no garantizan optimalidad y dependen fuertemente del diseño de sus movimientos. Una constante en ambas familias es que las reglas de descanso y de pago son no lineales y dependen del turno completo, por lo que la mayoría de los trabajos evalúan los turnos mediante una función procedimental y descartan los infactibles.

Entre las estrategias más prometedoras destacan: (i) los movimientos que operan sobre segmentos de turnos, como el intercambio de sufijos~\cite{DeLeone2011a} y la recombinación por asignación~\cite{Constantino2017}, que permiten fusionar turnos y reducir el tiempo ocioso; (ii) guiar la búsqueda con componentes específicos del costo, como el tiempo ocioso y el sobretiempo, en lugar del tiempo pagado total; y (iii) los híbridos entre métodos exactos y metaheurísticas, como LNS con reparación por \emph{branch and price}~\cite{Kletzander2026}, que representan la tendencia actual.

Como trabajo futuro, se propone explorar representaciones de asignación directa tramo--conductor con evaluación incremental del costo, que permiten movimientos baratos y un manejo explícito de las restricciones, ya sea descartando los vecinos infactibles o penalizándolos. También se propone estudiar la incorporación de la reprogramación dinámica ante eventos imprevistos, señalada como ventaja potencial por~\cite{Constantino2017}, y la integración con la etapa de \emph{crew rostering}. El modelo propuesto en la Sección~\ref{sec:modelo} constituye la base formal para la función de evaluación y el manejo de restricciones de la implementación de la siguiente entrega.

\section{Uso LLMs}
\begin{enumerate}
\item \textbf{Herramienta utilizada:} Claude Code (interfaz de línea de comandos de Anthropic), con el modelo Claude Opus 5.5.
\item \textbf{Propósito de uso:} lectura y síntesis de las fuentes entregadas por el curso (presentación del problema y artículos de De Leone et al. y Constantino et al.), búsqueda y verificación bibliográfica de referencias adicionales, redacción del borrador del informe, formulación del modelo matemático de la Sección~\ref{sec:modelo} y generación del código \LaTeX{} (incluida la figura).
\item \textbf{Prompts relevantes:} ``Necesito que me ayudes con la entrega 1 del proyecto de IA [\dots] dentro de las 3 carpetas repartidas está toda la info que necesitas. Necesito que hagas la entrega 1 completa. Después pasaremos a la revisión del mismo''. Las carpetas contenían las reglas del proyecto, la plantilla \LaTeX{}, la presentación del problema y los dos artículos fuente.
% REVISAR: actualizar los puntos 4 y 5 para que describan fielmente la revisión que realmente se hizo.
\item \textbf{Nivel de edición:} el borrador fue generado por la herramienta y luego revisado por el autor, quien verificó la coherencia del modelo y la redacción.
\item \textbf{Validación:} el contenido se contrastó con los artículos fuente~\cite{DeLeone2011a,Constantino2017} y con la presentación del problema entregada por el ayudante. Los datos bibliográficos de las referencias adicionales se verificaron mediante búsquedas en las páginas de las editoriales (AAAI, JAIR, Elsevier, Taylor \& Francis).
\end{enumerate}

\bibliographystyle{plain}
\bibliography{Referencias}

\end{document}
